\chapter{Text Data Analysis}
\label{chap:text}
% Text chapter, owner: Tô Anh Phát
% Numbers come from notebooks/text.ipynb (final local run); figures are exported from the same run with English labels
\sloppy

\section{Problem and Objectives}
Vietnamese e-commerce platforms contain many reviews that do not describe the product: copy-pasted praise,
reviews that only thank the shop or the courier, and text that is not a review at all.
Such reviews mislead buyers. This chapter predicts, from the review text (and the star rating),
whether a review is spam (\textbf{task 1}, two classes) and which type of spam it is (\textbf{task 2}, four classes).

Spam is the minority class (26\% of reviews), so accuracy would hide errors on the rare classes.
The main metric is therefore the \textbf{macro F1-score}, which is also the metric of the original paper,
so the results can be compared with its deep learning models.

\section{Dataset}
\begin{itemize}[leftmargin=1.5cm]
    \item Name: UIT-ViSpamReviews, version 1.
    \item URL: \url{https://github.com/sonlam1102/vispamdetection} (official zip from the authors).
    \item Source: reviews of recently best-selling products on major Vietnamese e-commerce platforms, at most 15 reviews per product.
          Three trained annotators labelled every review; the final label is the majority vote, with Cohen's kappa 0.72 (task 1) and 0.68 (task 2).
    \item Size: 19,870 reviews with 4 columns: \texttt{Rating} (1--5 stars), \texttt{Comment} (text), \texttt{Label} (task 1) and \texttt{SpamLabel} (task 2).
          The authors provide a fixed split of 14,306 train / 1,590 dev / 3,974 test reviews.
    \item License / citation: CC BY-NC 4.0, research use. C. Van Dinh, S. T. Luu, A. G.-T. Nguyen,
          \textit{Detecting Spam Reviews on Vietnamese E-commerce Websites}, ACIIDS 2022.
\end{itemize}

\begin{table}[H]
\centering
\caption{Spam types and their share of the training set}
\label{tab:text-labels}
\begin{tabular}{>{\raggedright\arraybackslash}p{2.6cm}>{\raggedright\arraybackslash}p{8.2cm}rr}
\toprule
Label & Definition (paper) & Train & Share\\
\midrule
NO-SPAM & Normal review & 10,585 & 74.0\%\\
SPAM-1 & Fake review: exaggerated praise or criticism, not about the real product & 195 & 1.4\%\\
SPAM-2 & Only talks about the shop, the brand or the delivery & 1,087 & 7.6\%\\
SPAM-3 & Not a review: advertising, meaningless characters, off-topic chat & 2,439 & 17.0\%\\
\bottomrule
\end{tabular}
\end{table}

\textbf{Leakage policy.} The three sets are loaded into one table only to check data quality and to find
duplicates across sets. Every exploratory analysis and every decision that uses labels is made on the
\textbf{training set only}; the dev set is used to compare models and the test set is used once, at the end.

\section{Data Quality Checks}
\begin{table}[H]
\centering
\caption{Data quality issues and how they were handled}
\label{tab:text-quality}
\begin{tabular}{>{\raggedright\arraybackslash}p{4.4cm}>{\raggedright\arraybackslash}p{4.2cm}>{\raggedright\arraybackslash}p{5.8cm}}
\toprule
Issue & Size & Handling\\
\midrule
Missing values, inconsistent labels & none & --\\
Same text repeated (ignoring case and spaces) & 121 texts (0.6\%) & Kept: mostly very short texts such as ``ok''\\
Same text with different labels & 7 texts (5 inside train, 11 rows) & The 5 conflicts found inside train are removed from train only\\
Test reviews copied from train & 41 (1.0\% of test) & Kept to stay comparable with the paper; impact measured on the test set\\
Text not in Unicode NFC form & 214 train reviews (1.5\%) & Normalised to NFC\\
Reviews empty after cleaning & 50 train reviews, all SPAM-3 & Kept; captured by character features and an \texttt{is\_empty} flag\\
Abbreviations and teencode (``ko'', ``sp'') & 23.6\% of train reviews & Kept; expanding them was tested and did not help (Table~\ref{tab:text-cv})\\
No Vietnamese diacritics & 678 train reviews (4.7\%) & Kept; 382 of them are gibberish, English or codes, 378 of which are SPAM-3\\
\bottomrule
\end{tabular}
\end{table}

\textbf{Text normalisation.} \texttt{normalize\_text} converts Unicode to NFC, lower-cases the text, replaces links with the token
\texttt{urltoken}, replaces punctuation and emoji with spaces and collapses spaces. This is the standard cleaning, but it removes
signals that matter here: the 50 reviews made only of emoji or symbols become empty, and all of them are SPAM-3.

\textbf{Word segmentation.} A Vietnamese word can have several syllables (\textit{sản phẩm} = ``product''), and spaces only
separate syllables. Words were segmented with \texttt{underthesea}, which reports about 97\% F1 on the VLSP 2013 news corpus.
Tests on real reviews show that it handles well-written text but merges words across a missing sentence break
(\textit{nhanh shop} becomes \textit{nhanh\_shop}) and only partly segments text without diacritics. The tool is therefore run
before punctuation is removed, since punctuation tells it where phrases end.

\textbf{A collection artefact.} The most frequent review length is exactly 300 characters: 262 training reviews end mid-word,
and 74\% of them are SPAM-3. The data was probably truncated when it was collected. Since new reviews will not be truncated,
this length is not used as a feature.

\section{Exploratory Data Analysis}
One review in four is spam, and SPAM-1 is 54 times rarer than NO-SPAM, which calls for macro F1 and class weights.
93.5\% of the reviews give 5 stars; 96\% of spam and 93\% of normal reviews are 5-star, so the rating carries little information.
The median review has 20 syllables (17 words); SPAM-1 reviews are the longest (median 50 syllables), while SPAM-3 ranges
from empty to long advertisements.

\subsection{Question 1: Which Writing Signals Separate the Classes?}
The cleaning step removes emoji, links, line breaks and repeated letters, and the vocabulary cannot see length or missing diacritics.
Ten simple signals were therefore measured for each class (Figure~\ref{fig:text-signals}).

\begin{figure}[H]
\centering
\figureplaceholder{text/text_signals.png}{Share of each class showing each writing signal.}
\caption{SPAM-3 stands out: 28.2\% have at most 3 syllables (2.1\% of NO-SPAM), 23.2\% have no diacritics (0.9\%) and 13.2\% repeat a letter at least 10 times (0.2\%). SPAM-1 is written on several lines (39.5\% against 14.9\%) and rarely uses abbreviations (9.2\% against 26.6\%). Digits are equally common in every class.}
\label{fig:text-signals}
\end{figure}

\subsection{Question 2: Does SPAM-2 Really Talk Only About the Shop and the Delivery?}
\begin{figure}[H]
\centering
\figureplaceholder{text/text_topics.png}{Share of each class mentioning each topic.}
\caption{69\% of SPAM-2 reviews mention delivery or packaging but only 30\% mention the product, and 21\% say they have not used the product yet (0--2\% in other classes). However, 31\% of normal reviews also only mention delivery or the shop, so the SPAM-2 / NO-SPAM boundary is blurred.}
\label{fig:text-topics}
\end{figure}

\subsection{Question 3: Are Fake Reviews Copy-Pasted Templates?}
A genuine review rarely shares 6 consecutive syllables with another review by chance. Phrases of 6 syllables that appear in
at least 10 different reviews were counted for each class.

\begin{figure}[H]
\centering
\figureplaceholder{text/text_templates.png}{Shared 6-syllable phrases per class.}
\caption{48.7\% of SPAM-1 reviews contain a shared 6-syllable phrase, against 11.3\% of normal reviews. Most top SPAM-1 phrases are consecutive pieces of one template (\textit{``... tôi không thể trông đợi gì hơn. Cảm ơn nhà cung cấp ...''}); 41 of the 195 fake reviews (21\%) paste one of two templates verbatim.}
\label{fig:text-templates}
\end{figure}

\textbf{Vocabulary.} The training set contains 12,307 distinct syllables but 20,412 distinct words, and 64\% of the words occur only once.
A log-odds comparison of word unigrams and bigrams confirms the topics above: \textit{chưa sử\_dụng} (``not used yet'') is the most
typical SPAM-2 term, NO-SPAM reviews describe concrete experience, and SPAM-3 terms come from unaccented advertising text messages.

\section{Data Preprocessing}
\subsection{Pipeline and Data Leakage Prevention}
All text transformations are row-wise and are applied to the three sets by one function, \texttt{model\_inputs}, which is also used
for new reviews. Every statistic (TF-IDF vocabulary and idf, scaler parameters, one-hot categories) is learned inside a
\texttt{Pipeline} with a \texttt{ColumnTransformer}, so it only comes from the training data, including inside cross-validation.
Dev and test keep their original size (1,590 and 3,974 reviews); train has 14,295 reviews after removing the 11 conflicting rows.

\subsection{Features}
\begin{table}[H]
\centering
\caption{Feature groups and the EDA finding that motivates each one}
\label{tab:text-features}
\begin{tabular}{>{\raggedright\arraybackslash}p{6.2cm}>{\raggedright\arraybackslash}p{8.2cm}}
\toprule
Features & Reason\\
\midrule
TF-IDF on syllable or word 1--2 grams (\texttt{sublinear\_tf}, \texttt{min\_df}=2) & Frequent words appear in every class (idf); spam repeats letters and phrases (sublinear tf); 64\% of words occur once (\texttt{min\_df})\\
TF-IDF on character 1--5 grams (\texttt{char\_wb}) of the text \emph{before} punctuation and emoji are removed & Keeps emoji, repeated letters, gibberish and unaccented text that cleaning erases\\
10 signals: at most 3 syllables, empty, several lines, no diacritics, emoji, link, letter repeated 4+ and 10+ times, abbreviation, $\log(1+\text{length})$; scaled with \texttt{MaxAbsScaler} & Figure~\ref{fig:text-signals}; length differs between classes and is right-skewed\\
Star rating, one-hot encoded & Weak according to the EDA; kept only to test it\\
\bottomrule
\end{tabular}
\end{table}

\section{Models and Experimental Setup}
The models are a majority-class baseline (\texttt{DummyClassifier}), Complement Naive Bayes (designed for imbalanced text),
logistic regression and a linear SVM, the last two with \texttt{class\_weight="balanced"}.
Deep models such as PhoBERT were not trained (they need a GPU); their scores are taken from the paper.
The selection was done in three steps:
\begin{enumerate}[leftmargin=1.5cm]
    \item Four text representations $\times$ three models, trained on train and scored on dev, for both tasks (Figure~\ref{fig:text-models}).
    \item The dev set has only 25 SPAM-1 reviews, so small differences on dev can be luck. The next choices (signals or not,
          rating or not, regularisation $C \in \{0.1, 0.3, 1, 3\}$, expanding abbreviations or not) were made with
          3-fold stratified cross-validation on the training set (Table~\ref{tab:text-cv}).
    \item The chosen pipeline was refitted on train + dev and evaluated once on the test set.
\end{enumerate}

\begin{figure}[H]
\centering
\figureplaceholder{text/text_model_comparison.png}{Macro F1 on dev for every representation and model.}
\caption{For task 1 the representation matters more than the model: with a linear SVM, syllables score 78.9, words 78.8, characters 81.7 and words + characters 83.1. Word segmentation adds at most 1.1 points over syllables (logistic regression), because syllable bigrams already rebuild most compound words; character features add 2.8 points. For task 2, syllables (73.9) and words + characters (73.3) are tied.}
\label{fig:text-models}
\end{figure}

\begin{table}[H]
\centering
\caption{Macro F1 (\%) with 3-fold cross-validation on train ($C = 0.3$, best representation of each task)}
\label{tab:text-cv}
\begin{tabular}{lrr}
\toprule
Setting & Task 1 & Task 2\\
\midrule
Text features only & 84.1 & 66.0\\
+ 10 signals & 84.4 & 68.6\\
+ 10 signals + rating (selected) & \textbf{84.4} & \textbf{68.7}\\
+ 10 signals + rating, abbreviations expanded & 84.4 & 69.0\\
\bottomrule
\end{tabular}
\end{table}

The signals add 2.6 points for task 2, whose model has no character features, but only 0.3 points for task 1, where the
character n-grams already capture most of them. The rating and the expansion of abbreviations change the score by at most
0.3 points, so abbreviations are kept as written. $C = 0.3$ is the best value for both tasks.
The final models are words + characters + signals + rating with a linear SVM for task 1, and syllables + signals + rating with
a linear SVM for task 2.

\section{Results and Evaluation}
\begin{table}[H]
\centering
\caption{Results on the test set (\%); the paper does not report the accuracy of TextCNN, LSTM and GRU}
\label{tab:text-results}
\begin{tabular}{lrrrr}
\toprule
& \multicolumn{2}{c}{Task 1} & \multicolumn{2}{c}{Task 2}\\
\cmidrule(lr){2-3}\cmidrule(lr){4-5}
Model & Accuracy & Macro F1 & Accuracy & Macro F1\\
\midrule
Baseline (always NO-SPAM) & 72.1 & 41.9 & 72.1 & 21.0\\
TextCNN / LSTM / GRU (paper) & -- & 77.2--77.9 & -- & 64.7--66.6\\
PhoBERT (paper) & 90.0 & \textbf{86.9} & 88.9 & 72.2\\
TF-IDF + signals + linear SVM (ours) & 88.1 & 85.4 & 86.1 & \textbf{74.0}\\
\bottomrule
\end{tabular}
\end{table}

The linear model reaches 85.4 macro F1 on task 1, 1.5 points below PhoBERT and 7.5--8.2 points above the paper's
recurrent and convolutional networks. On task 2 it scores 74.0, 1.8 points above PhoBERT, while training in a few seconds on a CPU.
The deep learning scores are those reported by the paper and were not reproduced here.
Removing the 41 test reviews that also appear in train changes macro F1 by at most 0.1 points (85.3 and 74.0),
so this overlap does not inflate the results.

\begin{figure}[H]
\centering
\figureplaceholder{text/text_confusion.png}{Confusion matrices on the test set.}
\caption{Task 1 detects 80.6\% of spam (precision 77.6\%). In task 2, SPAM-2 is the hardest class (F1 58.6): 123 SPAM-2 reviews (39\%) are predicted as NO-SPAM, the same main error as PhoBERT in the paper (177 reviews). SPAM-1 still reaches F1 68.4 with only 61 test reviews.}
\label{fig:text-confusion}
\end{figure}

\textbf{Error analysis.} For task 1, the share of correct predictions is 91.0\% for NO-SPAM, 88.0\% for SPAM-3, 70.5\% for SPAM-1
and 65.5\% for SPAM-2. Reading the wrong predictions shows four causes:
(1) SPAM-2 reviews that describe the delivery in detail and briefly mention the product, the blurred boundary of Figure~\ref{fig:text-topics};
(2) fake reviews that do not use a template and read like normal praise;
(3) doubtful labels, for example \textit{``Đánh giá cho shop 5 sao vì sự nhiệt tình của shop''} (``5 stars for the shop's kindness'')
labelled NO-SPAM although it only talks about the shop, in line with the annotators' kappa of 0.68--0.72;
(4) normal reviews written in a ``teen'' style with repeated letters, a signal the model learned from spam.

\begin{figure}[H]
\centering
\figureplaceholder{text/text_weights.png}{Largest word weights of the task 1 model.}
\caption{The model's weights agree with the EDA: \textit{chưa sử\_dụng} and \textit{chưa dùng} (``not used yet'') push most strongly towards spam, while concrete product comments such as \textit{hàng tốt} (``good goods'') and \textit{chất\_lượng tốt} (``good quality'') push towards NO-SPAM. Among the signals, length, letters repeated 10+ times, at most 3 syllables and links have the largest positive weights.}
\label{fig:text-weights}
\end{figure}

\section{Question-Driven Figure}
\textbf{Question:} Does word segmentation help more than keeping the information that text cleaning removes?

Figure~\ref{fig:text-models} and Table~\ref{tab:text-cv} answer it. Segmenting words changes macro F1 by at most 1.1 points compared
with plain syllables. Restoring what cleaning erases helps more: character n-grams on the uncleaned text add 2.8 points on task 1,
and the 10 hand-crafted signals add 2.6 points on task 2. For Vietnamese spam, \emph{how} a review is written is as informative as
the words it uses.

\section{Limitations and Future Work}
\begin{itemize}[leftmargin=1.5cm]
    \item Label noise: annotator agreement is 0.68--0.72 and 7 texts have conflicting labels; some errors in the error analysis are
          probably labelling errors rather than model errors.
    \item SPAM-1 is very small (195 training and 61 test reviews), so its F1 is unstable; part of the data is truncated at 300 characters.
    \item The abbreviation list and the signals were defined by hand; the word segmenter fails on unaccented text and on missing sentence breaks;
          a bag-of-words model cannot capture long context or sarcasm.
    \item Only the authors' single split was used, and the deep learning scores come from the paper rather than from the same pipeline.
    \item Future work: fine-tune PhoBERT on the same cleaned training data for a fair comparison and combine it with the hand-crafted signals,
          tune the task 1 decision threshold when catching more spam matters more than precision, and relabel the short conflicting reviews.
\end{itemize}

\section{References}
\begin{itemize}[leftmargin=1.5cm]
    \item C. Van Dinh, S. T. Luu, A. G.-T. Nguyen, \textit{Detecting Spam Reviews on Vietnamese E-commerce Websites},
          ACIIDS 2022, arXiv:2207.14636. Dataset: \url{https://github.com/sonlam1102/vispamdetection}, CC BY-NC 4.0.
    \item underthesea contributors, \textit{underthesea: Vietnamese NLP Toolkit}, \url{https://github.com/undertheseanlp/underthesea}.
    \item F. Pedregosa et al., \textit{Scikit-learn: Machine Learning in Python}, JMLR 12 (2011) 2825--2830.
\end{itemize}

\fussy
